\documentclass[10pt]{article}
\usepackage[margin=20mm]{geometry}
\usepackage{amsmath,amssymb,mathtools,fontspec,unicode-math,xurl,hyperref}
\IfFontExistsTF{DejaVu Serif}{\setmainfont{DejaVu Serif}}{\setmainfont{Latin Modern Roman}}
\setmathfont{latinmodern-math.otf}
\setlength{\parindent}{0pt}\setlength{\parskip}{6pt}\sloppy
\title{Initial review of
On the Serre conjecture for Artin characters
in the geometric case}
\author{}
\date{26 October 2026}
\begin{document}
\maketitle
Abstract: This review examines Tomoyuki Abe’s publisher version of record, Inventiones mathematicae 241 (2025), 605–622. Two local mathematical corrections are established: a fixed-point scheme need not equal its closed fiber, and the numerical discriminant composition formula needs a ramification factor. Both admit bounded repairs. Neither supplies a counterexample to the main theorem. The deeper external nearby-cycle and reduction theorems remain explicit dependencies of this review.\par
\section*{Identity, outcome and coverage}
M25-0084 is bound by its source manifest to the 18-page publisher PDF, DOI 10.1007/s00222-025-01345-w [paper]. A PDF page \(r\) corresponds to printed page \(604+r\). The outcome is Major: one core proof-step error, together with one Minor numerical-formula error. No counterexample to the main theorem is asserted. All 1,053 extracted text lines were read consecutively, including references. Fourteen original pages were actually viewed: PDF pages 1, 2 and 6–17. Exact symbolic calculations were run using independently written code. No author source or computational code is present in the assigned manifest; this absence is not a scientific deficiency.\par
\section*{Main claim, quantifiers and internal chain}
The input is a regular local ring \(A\), maximal ideal \(\mathfrak m\), and finite group \(G\), with \(A^G\) noetherian, \(A\) finite over \(A^G\), every \(A/I_\sigma\) of finite length for \(\sigma\ne1\), and the specified equality of residue fields. Here \(I_\sigma=(a-\sigma(a):a\in A)\). The character function is \[a_G(\sigma)=-\operatorname{length}_A(A/I_\sigma)\quad(\sigma\ne1),
 \qquad a_G(1)=-\sum_{\sigma\ne1}a_G(\sigma).\] Theorem on printed page 606 proves realization and \(\mathbb Q_\ell\)-rationality when \(A\) contains a finite field, for every prime \(\ell\) invertible in \(A\). It is a positive-characteristic statement. It does not establish the mixed-characteristic generalization or the separate \(p\)-adic conjecture of Section 3.2.\par
Theorem 1.3 gives base-change compatibility of the relative Lefschetz–Verdier trace for locally acyclic complexes. Lemma 1.4 descends bounded constructible derived objects through affine inverse limits. Theorem 1.5 obtains a universally locally acyclic extension after an alteration, retaining the specified endomorphism. These supply the relative trace family used in Theorem 1.6.\par
Theorem 1.6 assumes a flat finite-type map from a regular scheme to a strictly henselian smooth trait over algebraically closed \(k\), an action over the trait, an isolated fixed-point locus, and smoothness away from that point. It asserts, only for \(\sigma\ne1\), \[\operatorname{Tr}\!\left(\sigma;
 \Psi_{\mathrm{pr}'}\mathcal L(ft')_{(x,\infty')}\right)
       =-a_{X,G}(\sigma).\] The connected alteration of the parameter line identifies the two specializations of the relative trace. At zero the trace is the diagonal–graph intersection length; at infinity the nearby cycles are supported at the isolated point. Proposition 2.5 identifies their virtual class with minus the total-dimension class of vanishing cycles. Section 2.6 then invokes the Kato–Saito–Saito reduction, paper reference [13], Lemma 5.3.\par
\section*{Problem 1: support at a point is not schematic containment in its fiber}
\subsection*{Location}
In Theorem 1.6’s proof, printed page 613 (PDF page 9), the display immediately before the geometric Claim starts with \[(X_i)^\sigma\simeq (X_i)^\sigma\times_{S_i}\{s_i\}.\] The preceding reduction establishes that the fixed scheme has underlying support a single point. It does not establish that the base uniformizer annihilates its structure sheaf.\par
\subsection*{A full-hypothesis counterexample to the displayed isomorphism}
Let \(k\) be algebraically closed of characteristic \(2\), take \(\ell=3\), and consider \[X_i=\operatorname{Spec}k[t,u]/(u^2+t^2u+t),\qquad
 S_i=\mathbb A^1_t,\qquad f(t,u)=t.\] Let \(G=C_2=\{1,\sigma\}\) act by \[\sigma(t)=t,\qquad \sigma(u)=u+t^2.\] The defining polynomial is invariant, since the two \(t^4\) terms cancel in characteristic \(2\). Applying \(\sigma\) twice is the identity, and \(\sigma\ne1\). The map \(f\) is finite flat of degree two because the polynomial is monic in \(u\). The curve is smooth over \(k\): \(\partial(u^2+t^2u+t)/\partial t=1\). Away from \(t=0\), its derivative in \(u\), namely \(t^2\), is invertible, so \(f\) is etale there. The fiber at zero has the single underlying point \(x=(0,0)\).\par
The fixed ideal is exactly \((t^2)\). Indeed the difference of the generator \(u\) is \(t^2\), and all polynomial differences are multiples of it. At \(x\) its coordinate ring is \[k[t,u]/(u^2+t^2u+t,t^2)
       \simeq k[u]/(u^4),\qquad t\longmapsto u^2.\] It has length four and \(t\ne0\) in it. Taking the closed fiber instead gives \[k[t,u]/(u^2+t^2u+t,t)
       \simeq k[u]/(u^2),\] of length two. Thus the displayed isomorphism fails. The fixed locus has precisely the required singleton underlying support. After strict henselization at \(t=0\), the example meets the local assumptions of Theorem 1.6 as well. The group action, regularity, flatness and smooth-away-from-\(x\) conditions have all been checked, rather than inferred from the support alone.\par
\subsection*{Counterevidence, impact and repair}
The required conclusion \(a_{X_i,G}=a_{X,G}\) remains true. For a finite-length algebra \(B\) supported at a \(k\)-rational closed point of \(S_i\), some power of its maximal ideal acts by zero. The local ring and its strict henselization have the same quotient by this power, since \(k\) is already separably closed. Hence \[B\otimes_{\mathcal O_{S_i,s_i}}\mathcal O^{\rm sh}_{S_i,s_i}
 \simeq B.\] Equivalently, a local Artinian ring with separably closed residue field is already strictly henselian. Fixed schemes commute with this base change because they are equalizers. Their entire nilpotent thickness, and therefore their length, is preserved.\par
Replace the false closed-fiber chain by this Artinian base-change argument. The correction does not require changing the theorem’s hypotheses. In the example the correct preserved length is four, not the closed-fiber length two. Later use of topological invariance of etale cohomology on the reduced support is a different operation and does not justify deleting thickness in an intersection length. The displayed reduction is an invalid derivation in the core trace theorem: it identifies the fixed-scheme lengths needed to pass from the finite-type geometric statement to the strictly henselian setting. The counterexample changes that length from four to two. Under claim-level calibration, this defect is Major. The Artinian base-change argument above repairs the reduction without changing the theorem’s hypotheses; no counterexample to the trace theorem or the final Artin-character theorem is asserted.\par
\section*{Problem 2: a missing weight in the numerical discriminant formula}
\subsection*{Location and interpretation}
Section 3.3, printed page 621 (PDF page 17), writes \[\operatorname{disc}(q)+\operatorname{disc}(g)
       =\operatorname{disc}(f),\qquad
 \mu_f-a_{X,G}(1)=\operatorname{disc}(g),
 \quad f=g\circ q.\] For ordinary local numerical discriminant/different exponents over an algebraically closed residue field, the second contribution must be multiplied by the ramification index of \(q\). If the notation for \(\operatorname{disc}(g)\) was intended to include a pullback or this weight, that convention must be stated: the numerical formula otherwise contradicts the following example.\par
\subsection*{Exact tame example}
Let \(k\) be algebraically closed of characteristic \(5\), choose \(\ell=7\), and work at zero on the curves \[X=\mathbb A^1_u,\qquad X/G=\mathbb A^1_v,\qquad
 S=\mathbb A^1_t,\qquad
 q:u\mapsto v=u^2,\quad g:v\mapsto t=v^3.\] The group \(C_2\) acts by \(u\mapsto-u\), and \(f:u\mapsto u^6\). The action is over \(S\), its nonidentity element fixes only zero, and \(f\) is finite flat and smooth away from zero. All rings are regular. One may make the same strict henselization as in Theorem 1.6.\par
The coefficients \(2,3,6\) are nonzero in \(k\). Therefore \[\mu_f=\operatorname{ord}_u(6u^5)=5,\qquad
 \operatorname{disc}(q)=1,\qquad \operatorname{disc}(g)=2.\] Also \(I_\sigma=(2u)=(u)\), so \(a_G(\sigma)=-1\) and \(a_G(1)=1\). The printed formula would give \(5=1+2\), or \(5-1=2\). Both are false. The correct weighted calculation is \(5=1+2\cdot2\).\par
More generally, for \(t=g(v)\), \(v=q(u)\), the chain rule gives \[\operatorname{ord}_u\frac{dt}{du}
 =\operatorname{ord}_u\frac{dv}{du}
   + e(q)\operatorname{ord}_v\frac{dt}{dv}.\] The valuation factor comes from substitution: a uniformizer \(v\) has \(u\)-valuation \(e(q)\). For the finite local quotient under the stated residue-field conditions, \(e(q)=|G|\). Thus the correctly normalized scalar statement is \[\mu_f-a_{X,G}(1)=|G|\,\operatorname{disc}(g).\]\par
\subsection*{Counterevidence and repair limits}
The final implication in Section 3.3 survives: when \(g\) is etale, its exponent vanishes, so equality at \(\sigma=1\) does hold. Since \(e(q)>0\), vanishing is unchanged by the factor. Theorem 1.6 itself expressly excludes \(\sigma=1\). The example’s Artin character is the sign representation of \(C_2\), which is rational and satisfies the main theorem.\par
Insert the factor, or rewrite the different identity as a divisor identity with the pullback \(q^*\) displayed. This repair is confined to the one-dimensional, finite generically separable local setting. No higher-dimensional conductor formula or extension to \(\sigma=1\) for arbitrary \(f\) follows from it. The severity is Minor.\par
\section*{Checks that did not establish additional defects}
Conjugation carries \(I_\sigma\) to the corresponding conjugate ideal, so the length function is a class function. Its definition forces \(\sum_{\sigma\in G}a_G(\sigma)=0\). If it is a genuine character, this excludes a trivial summand; it does not itself prove nonnegative multiplicities. For \(G=\{1\}\), the function is zero, realized by the zero representation. In dimension zero the residue-field assumption forces a trivial action and gives the augmentation character. These boundary checks cause no contradiction.\par
The signs in Theorem 1.6 and Proposition 2.5 cancel correctly: \(\operatorname{dt}(\Phi)=-[\Psi]\) and \(\operatorname{Tr}(\sigma;\Psi)=-a_G(\sigma)\), so the total-dimension trace is \(a_G(\sigma)\) for \(\sigma\ne1\). The proof uses support away from \(x\) through Proposition 2.5(1); that support calculation is reproduced from independent nearby-cycle and Fourier-vanishing inputs, not deduced from the trace theorem being proved.\par
For the local diagonal–graph intersection at a smooth point, the graph and diagonal are regular embeddings of the expected codimension. The isolated fixed ideal is a system of parameters in a regular local ring, hence a regular sequence. This explains why the local intersection multiplicity is the scheme length rather than merely the reduced number of points.\par
Sections 2.2–2.3 deliberately choose compatible noncanonical Swan modules along a tower and then prove independence at the level of \(R_L(G)\). No canonical functor independent of all choices is needed for that conclusion. Projectivity of the Swan lattice gives exactness; torsion differences between lattices disappear after rationalization. The strict group-action model in Section 2.4 also permits alternating-trace additivity via an actual mapping cone; the general nonadditivity warning for arbitrary derived endomorphisms is not a counterexample to this construction. The derived tensor identities in Proposition 2.5 retain the \(\Lambda[\Delta]\)-perfectness requirement; replacing it by underlying vector-space perfectness would not be justified.\par
\section*{Unverified dependencies and evidence}
Theorem 1.5’s constructibility, cohomological-properness and universal-local-acyclicity steps were traced, including the finite descent of the complex and its endomorphism. The full Orgogozo theorems, the Lu–Zheng two-category construction, and Laumon’s local Fourier calculation were not independently re-proved. The Kato–Saito–Saito Lemma 5.3 reduction remains essential for the passage from geometric virtual trace identities to the full local-ring statement and positivity. A focused primary-source search identified the paper, but full access through JSTOR failed; this report does not claim to have re-verified that lemma. These are coverage limits, not asserted defects.\par
The official arXiv record checked on 26 October 2026 lists v2 dated 13 June 2025. It is contextual evidence only; this review is bound to the acquired publisher PDF, not to inferred equivalence with the author version. A targeted correction search did not identify a correction; the search was not exhaustive.\par
The independent Python/SymPy script performs exact finite-field polynomial differentiation, verifies the involution, and computes the two quotient Groebner bases. Its source and output are saved in the evidence directory. The local mathematical proofs above supply the hypotheses that symbolic output alone cannot certify. No author code was executed.\par
The report’s exact current TeX is compiled with authorized Tectonic 0.17.0. Build records contain actual logs, source/compiler/ dependency hashes before and after the final compile, and the resulting PDF hash. Every final report page is actually viewed and separately recorded. Coverage, findings and provenance are hash-bound in the per-paper delivery state. `human\_reviews=[]` and `scientific\_decision=pending`; initial delivery does not constitute scientific acceptance.\par
Focused independent acceptance is identified by token:\par
The current executor used a bounded read-only native helper for the existing problem, with actual roles and evidence recorded in provenance. The retained initial review is distinguished from this focused check; the scientific decision remains pending.\par
9 T. Abe, On the Serre conjecture for Artin characters in the geometric case, Invent. Math. 241 (2025), 605–622. https://doi.org/10.1007/s00222-025-01345-w. Assigned publisher PDF SHA-256: `f789468ee535bf3347550763fb4f9fce85efa8628247b01a9df2b04b008afb83`. T. Abe, official version record, arXiv:2405.19601, https://arxiv.org/abs/2405.19601, accessed 26 October 2026. K. Kato, S. Saito and T. Saito, Artin characters for algebraic surfaces, Amer. J. Math. 110 (1988), 49–75, paper reference [13]; especially Lemma 5.3. Full lemma not independently retrieved in this review.\par
\end{document}
